% Blueprint: the argument principle on a rectangle

\chapter{The argument principle on a rectangle}
\label{cha:argument_principle}

The exact Riemann--von Mangoldt formula of \Cref{cha:n_formula} rests on counting
the zeros of an entire function inside a rectangle by a boundary integral of its
logarithmic derivative. This is the rectangle \emph{argument principle}
(\Cref{thm:nf_argument_principle}), proved in this chapter.
\texttt{PrimeNumberTheoremAnd} supplies a complete \emph{single}-pole rectangle
residue calculus (\Cref{thm:nf_rectangle}); the work here is to (i) reduce a
finite set of zeros to the single-zero case by axis-parallel subdivision, and
(ii) identify the boundary contour integral with the boundary argument variation
of \Cref{cha:arg_variation}.

Throughout, $f:\mathbb C\to\mathbb C$, and $z,w\in\mathbb C$ are the corners of an
axis-aligned rectangle, normalised by the \emph{orientation hypothesis}
$z.\mathrm{re}\le w.\mathrm{re}$ and $z.\mathrm{im}\le w.\mathrm{im}$. We write
$R:=\texttt{Rectangle}\ z\ w$ for the closed rectangle
$[z.\mathrm{re},w.\mathrm{re}]\times[z.\mathrm{im},w.\mathrm{im}]$,
$\partial R:=\texttt{RectangleBorder}\ z\ w$ for its boundary, and
\[
  \operatorname{logDeriv} f \;=\; \frac{f'}{f}
  \;=\; \bigl(z\mapsto \tfrac{\operatorname{deriv} f\,z}{f\,z}\bigr)
\]
(\texttt{logDeriv}) for the logarithmic derivative. We abbreviate the integer
multiplicity of a zero by
\[
  m_f(\rho) \;:=\; \bigl(\texttt{meromorphicOrderAt}\ f\ \rho\bigr).\texttt{untopD}\ 0
  \;\in\;\mathbb Z ,
\]
which at an interior zero equals the (positive) order of vanishing; this is the
same quantity summed in \texttt{riemannZeta.N} via \Cref{lem:nf_N_eq_xi_count}.

\section{Finiteness of the zero set in the rectangle}
\label{sec:ap_finite}

\begin{lemma}[Zeros in a rectangle are finite]
\label{lem:ap_finite_zeros}
\lean{finite_rectangle_inter_setOf_eq_zero}
\leanok
Let $U\subseteq\mathbb C$ be open and connected, let $f$ be analytic on $U$
($\texttt{AnalyticOnNhd}\ \mathbb C\ f\ U$) and not identically zero on $U$, and let
$R=\texttt{Rectangle}\ z\ w\subseteq U$. Then the zero set
\[
  Z \;:=\; \{\rho\in R : f(\rho)=0\}
\]
is finite.
\end{lemma}

\begin{proof}
\leanok
On a connected open set on which an analytic function is not identically zero, the
zero set is \emph{codiscrete}: $f^{-1}\{0\}\in\texttt{codiscreteWithin}\ U$
(\texttt{AnalyticOnNhd.preimage\_mem\_codiscreteWithin}, whose hypothesis ``$f$ not
locally zero'' is supplied by the identity theorem
\texttt{AnalyticOnNhd.eqOn\_zero\_of\_preconnected\_of\_frequently\_eq\_zero} from
connectedness and $f\not\equiv0$). Concretely, every zero of $f$ in $U$ is isolated:
$\texttt{analyticOrderAt}\ f\ \rho$ is finite (not $\top$) at each $\rho\in U$, so by
\texttt{AnalyticAt.eventually\_eq\_zero\_or\_eventually\_ne\_zero} the zeros have no
accumulation point in $U$. Since $R$ is compact
(\texttt{IsCompact.reProdIm isCompact\_uIcc isCompact\_uIcc}) and $Z\subseteq R\subseteq U$
has no accumulation point, $Z$ is finite: an infinite subset of the compact set $R$
would have an accumulation point lying in $R\subseteq U$, contradicting isolation.
(This is the rectangle analogue of \texttt{riemannZeta.zeroes\_on\_Compact\_finite}
in \texttt{PrimeNumberTheoremAnd}, proved by the same codiscreteness-plus-compactness
argument.)
\end{proof}

\section{The residue sum of the logarithmic derivative}
\label{sec:ap_residue_sum}

The per-zero residue is already available.

\begin{lemma}[Residue of the logarithmic derivative]
\label{lem:ap_logderiv_residue}
\lean{meromorphicOrderAt_eq_residue_logDeriv}
\leanok
Let $f$ be analytic at $\rho$ and not eventually zero there, of order
$n=\texttt{meromorphicOrderAt}\ f\ \rho\ge0$. Then there are $n\in\mathbb N$ and a
function $g$ analytic at $\rho$ with $g(\rho)\neq0$ such that $m_f(\rho)=n$,
$\operatorname{logDeriv} g$ is analytic at $\rho$, and
\[
  \operatorname{logDeriv} f \;=\; \frac{n}{\,z-\rho\,} + \operatorname{logDeriv} g
  \qquad\text{on a punctured neighbourhood of }\rho .
\]
In particular $\operatorname{logDeriv} f - \dfrac{m_f(\rho)}{\,z-\rho\,}$ extends
analytically across $\rho$, hence is $O(1)$ on $\mathring{\mathbb C}\!\setminus\!\{\rho\}$
near $\rho$ (the residue hypothesis of \Cref{thm:nf_rectangle}).
\end{lemma}

\begin{proof}\leanok
This is \texttt{meromorphicOrderAt\_eq\_residue\_logDeriv} in
\texttt{Zeta/Analysis/Complex/ArgumentPrinciple.lean}. The boundedness clause
is immediate: the difference equals $\operatorname{logDeriv} g$ near $\rho$, which is
analytic at $\rho$ and therefore continuous, hence bounded on a punctured
neighbourhood; equivalently $\bigl(\operatorname{logDeriv} f-m_f(\rho)/(z-\rho)\bigr)
=O_{\,\mathcal N[\neq]\rho}(1)$.
\end{proof}

\begin{lemma}[At most one zero: single residue]
\label{lem:ap_one_zero}
\lean{rectangleIntegral_logDeriv_eq_two_pi_I_of_subsingleton}
\leanok
\uses{thm:nf_rectangle, lem:ap_logderiv_residue}
Suppose $z.\mathrm{re}\le w.\mathrm{re}$, $z.\mathrm{im}\le w.\mathrm{im}$, that $f$
is analytic on a neighbourhood of $R=\texttt{Rectangle}\ z\ w$, that $f\neq0$ on
$\partial R$, and that $Z=\{\rho\in R:f(\rho)=0\}$ has at most one element. Then
\[
  \texttt{RectangleIntegral}\ (\operatorname{logDeriv} f)\ z\ w
  \;=\; 2\pi i \cdot \sum_{\rho\in Z} m_f(\rho).
\]
\end{lemma}

\begin{proof}
\leanok
\uses{thm:nf_rectangle, lem:ap_logderiv_residue}
If $Z=\varnothing$ then $f\neq0$ on all of $R$, so $\operatorname{logDeriv} f$ is
analytic on $R$ and \texttt{HolomorphicOn.vanishesOnRectangle} gives
$\texttt{RectangleIntegral}\ (\operatorname{logDeriv} f)\ z\ w=0$, matching the empty
sum.

If $Z=\{\rho\}$, then $\rho$ lies in the \emph{interior} of $R$: it is in $R$ and not
on $\partial R$ (where $f\neq0$), so $R\in\mathcal N(\rho)$
(\texttt{rectangle\_mem\_nhds\_iff}). On $R\setminus\{\rho\}$ we have $f\neq0$, so
$\operatorname{logDeriv} f$ is holomorphic there. By \Cref{lem:ap_logderiv_residue}
with $A:=m_f(\rho)$, the difference $\operatorname{logDeriv} f-A/(z-\rho)$ is $O(1)$
near $\rho$. Hence \texttt{ResidueTheoremOnRectangleWithSimplePole'} applies (with the
orientation hypotheses and $R\in\mathcal N(\rho)$) and yields
$\texttt{RectangleIntegral'}\ (\operatorname{logDeriv} f)\ z\ w = A$, i.e.
\[
  \texttt{RectangleIntegral}\ (\operatorname{logDeriv} f)\ z\ w
  = 2\pi i\,A = 2\pi i\,m_f(\rho),
\]
since $\texttt{RectangleIntegral'}=(2\pi i)^{-1}\texttt{RectangleIntegral}$. The sum
over $Z=\{\rho\}$ is exactly $m_f(\rho)$.
\end{proof}

\begin{theorem}[Rectangle residue sum for the logarithmic derivative]
\label{thm:ap_residue_sum}
\lean{rectangleIntegral_logDeriv_eq_two_pi_I_sum}
\leanok
\uses{lem:ap_finite_zeros, lem:ap_one_zero, thm:nf_rectangle}
Let $U$ be open and connected, $f$ analytic on $U$ and not identically zero, and
$R=\texttt{Rectangle}\ z\ w\subseteq U$ with $z.\mathrm{re}\le w.\mathrm{re}$,
$z.\mathrm{im}\le w.\mathrm{im}$ and $f\neq0$ on $\partial R$. Then, with
$Z=\{\rho\in R:f(\rho)=0\}$ (finite by \Cref{lem:ap_finite_zeros}),
\[
  \texttt{RectangleIntegral}\ (\operatorname{logDeriv} f)\ z\ w
  \;=\; 2\pi i \cdot \sum_{\rho\in Z} m_f(\rho).
\]
\end{theorem}

\begin{proof}
\leanok
\uses{lem:ap_finite_zeros, lem:ap_one_zero, thm:nf_rectangle}
Strong induction on the cardinality $\lvert Z\rvert$ (finite by
\Cref{lem:ap_finite_zeros}). If $\lvert Z\rvert\le1$ this is \Cref{lem:ap_one_zero}.

Assume $\lvert Z\rvert\ge2$ and pick two distinct $\rho_1,\rho_2\in Z$. Then either
their real parts differ or their imaginary parts differ. Consider the former case
(the latter is symmetric, using a horizontal cut). The finitely many real parts
$\{\rho.\mathrm{re}:\rho\in Z\}$ do not exhaust the open interval strictly between
$\min$ and $\max$ of two distinct values among them, so there is a real
$a\in(z.\mathrm{re},w.\mathrm{re})$ with $a\neq\rho.\mathrm{re}$ for every
$\rho\in Z$, and with at least one $\rho\in Z$ on each side $\rho.\mathrm{re}<a$,
$\rho.\mathrm{re}>a$ (choose $a$ strictly between two distinct real parts of $Z$
while avoiding the finite set of real parts). Split $R$ along the vertical segment
$\mathrm{re}=a$ into
\[
  R_1=\texttt{Rectangle}\ z\ (a+w.\mathrm{im}\,i),
  \qquad
  R_2=\texttt{Rectangle}\ (a+z.\mathrm{im}\,i)\ w .
\]
On $\partial R_1$ and $\partial R_2$ we have $f\neq0$: their borders consist of
pieces of $\partial R$ (where $f\neq0$) together with the new cut segment
$\{a+t i: t\in[z.\mathrm{im},w.\mathrm{im}]\}$, which contains no zero of $f$ because
$a$ avoids all real parts of $Z$. Consequently $\operatorname{logDeriv} f$ is
continuous on each border, hence \texttt{RectangleBorderIntegrable}
(\texttt{ContinuousOn.rectangleBorderIntegrable}), and the horizontal-split lemma
\texttt{RectangleIntegralHSplit'} (applicable since
$a\in[[z.\mathrm{re},w.\mathrm{re}]]$ and $f$ is border-integrable on $R$) gives
\[
  \texttt{RectangleIntegral}\ (\operatorname{logDeriv} f)\ z\ w
  = \texttt{RectangleIntegral}\ (\operatorname{logDeriv} f)\ z\ (a+w.\mathrm{im}\,i)
  + \texttt{RectangleIntegral}\ (\operatorname{logDeriv} f)\ (a+z.\mathrm{im}\,i)\ w .
\]
The zero sets of $f$ in $R_1$ and $R_2$ are $Z_1=\{\rho\in Z:\rho.\mathrm{re}<a\}$ and
$Z_2=\{\rho\in Z:\rho.\mathrm{re}>a\}$ (no zero has $\mathrm{re}=a$); these partition
$Z$ and both are proper nonempty subsets, so
$\lvert Z_1\rvert,\lvert Z_2\rvert<\lvert Z\rvert$. Apply the induction hypothesis to
$R_1\subseteq U$ and $R_2\subseteq U$ (each still satisfies the orientation and
border-nonvanishing hypotheses):
\[
  \texttt{RectangleIntegral}\ (\operatorname{logDeriv} f)\ z\ (a+w.\mathrm{im}\,i)
  = 2\pi i\sum_{\rho\in Z_1} m_f(\rho),
  \qquad
  \cdots = 2\pi i\sum_{\rho\in Z_2} m_f(\rho).
\]
Adding and using $\sum_{Z_1}+\sum_{Z_2}=\sum_{Z}$ (disjoint union) gives the claim.
\end{proof}

\section{The contour integral as boundary argument variation}
\label{sec:ap_argvar}

We now realise the boundary contour as a closed concatenation of four affine
segments and connect its integral of $\operatorname{logDeriv} f$ to the argument
variation of \Cref{cha:arg_variation}. Recall $\texttt{lineSeg}\ P\ Q\ (u)=P+u(Q-P)$
from \texttt{Zeta/Theta.lean}, and write the four corners of $R$
counterclockwise from the bottom-left as
\[
  P_0=z.\mathrm{re}+z.\mathrm{im}\,i,\quad
  P_1=w.\mathrm{re}+z.\mathrm{im}\,i,\quad
  P_2=w.\mathrm{re}+w.\mathrm{im}\,i,\quad
  P_3=z.\mathrm{re}+w.\mathrm{im}\,i .
\]

\begin{definition}[Boundary argument variation of a rectangle]
\label{def:ap_rect_argvar}
\lean{rectangleArgVariation}
\leanok
\uses{def:av_argVariation}
For $f$ continuous and nonvanishing on $\partial R$, define the counterclockwise
argument variation of $f$ around $\partial R$ as the sum over the four oriented
edges
\[
  \Delta_{\partial R}\arg f \;:=\;
  \sum_{k=0}^{3}\ \texttt{argVariation}\ 0\ 1\ \bigl(u\mapsto f(\texttt{lineSeg}\
  P_k\ P_{k+1}\ u)\bigr),
  \qquad P_4:=P_0 .
\]
The four summands are the variations along bottom ($P_0\to P_1$), right
($P_1\to P_2$), top ($P_2\to P_3$) and left ($P_3\to P_0$). This is the object
denoted $\Delta_{\partial R}\arg f$ in \Cref{cha:n_formula}.
\end{definition}

\begin{lemma}[Real part of the logarithmic-derivative integral]
\label{lem:ap_re_integral}
\lean{re_integral_div_eq_log_norm_sub}
\leanok
\uses{def:av_isContLog, lem:av_exists}
Let $\varphi:[a,b]\to\mathbb C$ be continuous, nowhere zero, and differentiable on
$(a,b)$ with derivative $\varphi'$ such that $\varphi'/\varphi$ is continuous on
$[a,b]$. Then
\[
  \operatorname{Re}\!\int_a^b \frac{\varphi'(u)}{\varphi(u)}\,du
  \;=\; \log\lVert \varphi(b)\rVert - \log\lVert \varphi(a)\rVert .
\]
\end{lemma}

\begin{proof}
\leanok
\uses{def:av_isContLog, lem:av_exists, lem:av_hasDerivAt}
This is the real-part companion to \Cref{lem:av_integral}, whose Lean statement
\texttt{argVariation\_eq\_im\_integral} concerns the imaginary part. Let $g$ be the anchored continuous logarithm of $\varphi$
furnished by \Cref{lem:av_exists}; by \Cref{lem:av_hasDerivAt} $g$ is a $C^1$
primitive of $\varphi'/\varphi$ on $(a,b)$, so the fundamental theorem of calculus
(\texttt{intervalIntegral.integral\_eq\_sub\_of\_hasDeriv\_right\_of\_le}, exactly as
in the proof of \texttt{argVariation\_eq\_im\_integral}) gives
$\int_a^b \varphi'/\varphi = g(b)-g(a)$. Taking real parts and using
$\exp(g(u))=\varphi(u)\Rightarrow\operatorname{Re} g(u)=\log\lVert\varphi(u)\rVert$
(\texttt{Complex.norm\_exp} / \texttt{Complex.exp\_re}: $\lVert\exp w\rVert
=\exp(\operatorname{Re} w)$, hence $\operatorname{Re} w=\log\lVert\exp w\rVert$) yields
$\operatorname{Re}\int_a^b\varphi'/\varphi=\log\lVert\varphi(b)\rVert
-\log\lVert\varphi(a)\rVert$.
\end{proof}

\begin{lemma}[Side integral: imaginary part is the $\arg$ variation, real part is $\log\lVert\cdot\rVert$]
\label{lem:ap_side_integral}
\lean{argVariation_lineSeg_eq_im_integral}
\leanok
\uses{def:av_argVariation, lem:av_integral, lem:ap_re_integral}
Let $P,Q\in\mathbb C$ and let $f$ be holomorphic and nonvanishing on a neighbourhood
of the segment $\texttt{lineSeg}\ P\ Q\,([0,1])$. Then $u\mapsto f(\texttt{lineSeg}\
P\ Q\ u)$ is nowhere zero and, writing
$J:=\int_0^1 \operatorname{logDeriv} f(\texttt{lineSeg}\ P\ Q\ u)\,(Q-P)\,du$,
\[
  \texttt{argVariation}\ 0\ 1\ \bigl(f\circ\texttt{lineSeg}\ P\ Q\bigr)
  \;=\; \operatorname{Im} J ,
  \qquad
  \operatorname{Re} J \;=\; \log\lVert f(Q)\rVert-\log\lVert f(P)\rVert .
\]
\end{lemma}

\begin{proof}
\leanok
\uses{lem:av_integral, lem:ap_re_integral}
The path $\gamma(u)=\texttt{lineSeg}\ P\ Q\ u=P+u(Q-P)$ is affine, hence smooth with
constant derivative $\gamma'(u)=Q-P$, and its image is the segment. Set
$\varphi:=f\circ\gamma$; it is continuous, nowhere zero (on a neighbourhood of the
image $f\neq0$), and $C^1$ with $\varphi'(u)=f'(\gamma(u))\,\gamma'(u)
=\operatorname{logDeriv} f(\gamma(u))\,\varphi(u)\,(Q-P)$, so $\varphi'/\varphi
=\operatorname{logDeriv} f(\gamma(u))\,(Q-P)$ is continuous; thus $J=\int_0^1
\varphi'/\varphi$. The imaginary part is \Cref{lem:av_integral}
($\operatorname{logDeriv} f=f'/f$, $\gamma'=Q-P$):
\[
  \texttt{argVariation}\ 0\ 1\ (f\circ\gamma)
  = \operatorname{Im}\!\int_0^1 \frac{f'(\gamma(u))\,\gamma'(u)}{f(\gamma(u))}\,du
  = \operatorname{Im} J .
\]
The real part is \Cref{lem:ap_re_integral} applied to $\varphi$:
$\operatorname{Re} J = \log\lVert\varphi(1)\rVert-\log\lVert\varphi(0)\rVert
= \log\lVert f(Q)\rVert-\log\lVert f(P)\rVert$.
\end{proof}

\begin{lemma}[Contour integral is $i$ times the boundary argument variation]
\label{lem:ap_integral_eq_argvar}
\lean{rectangleIntegral_logDeriv_eq_I_mul_rectangleArgVariation}
\leanok
\uses{def:ap_rect_argvar, lem:ap_side_integral, thm:nf_rectangle}
Suppose $f$ is holomorphic on a neighbourhood of $R=\texttt{Rectangle}\ z\ w$ and
$f\neq0$ on $\partial R$. Then
\[
  \texttt{RectangleIntegral}\ (\operatorname{logDeriv} f)\ z\ w
  \;=\; i\cdot \Delta_{\partial R}\arg f .
\]
\end{lemma}

\begin{proof}
\leanok
\uses{lem:ap_side_integral, def:ap_rect_argvar}
By definition
$\texttt{RectangleIntegral}\ g\ z\ w
=\texttt{HIntegral}\ g\ z.\mathrm{re}\ w.\mathrm{re}\ z.\mathrm{im}
-\texttt{HIntegral}\ g\ z.\mathrm{re}\ w.\mathrm{re}\ w.\mathrm{im}
+\texttt{VIntegral}\ g\ w.\mathrm{re}\ z.\mathrm{im}\ w.\mathrm{im}
-\texttt{VIntegral}\ g\ z.\mathrm{re}\ z.\mathrm{im}\ w.\mathrm{im}$,
where $\texttt{HIntegral}\ g\ x_1\ x_2\ y=\int_{x_1}^{x_2} g(x+yi)\,dx$ and
$\texttt{VIntegral}\ g\ x\ y_1\ y_2=i\int_{y_1}^{y_2} g(x+yi)\,dy$. Reparametrising
each affine edge by $u\in[0,1]$ ($x=z.\mathrm{re}+u(w.\mathrm{re}-z.\mathrm{re})$ on
horizontal edges, $y=z.\mathrm{im}+u(w.\mathrm{im}-z.\mathrm{im})$ on vertical edges,
the linear change of variables contributing exactly the factor $Q-P=\gamma'$), the
four signed terms are precisely the four oriented edge integrals
$J_k:=\int_0^1 \operatorname{logDeriv} f(\texttt{lineSeg}\ P_k\ P_{k+1}\ u)\,(P_{k+1}-P_k)\,du$:
\[
  \texttt{RectangleIntegral}\ (\operatorname{logDeriv} f)\ z\ w
  = J_0 + J_1 + J_2 + J_3 ,
\]
the signs of \texttt{RectangleIntegral} matching the counterclockwise orientation
(bottom $P_0\to P_1$ as $+\texttt{HIntegral}$ at $z.\mathrm{im}$; top $P_2\to P_3$ as
$-\texttt{HIntegral}$ at $w.\mathrm{im}$, i.e. right-to-left; right $P_1\to P_2$ as
$+\texttt{VIntegral}$ at $w.\mathrm{re}$; left $P_3\to P_0$ as $-\texttt{VIntegral}$
at $z.\mathrm{re}$, i.e. top-to-bottom). Each closed edge is a compact subset of
$\partial R$, on which $f\neq0$; since $\{z:f(z)\neq0\}$ is open (continuity of $f$)
and contains the edge, it contains an open neighbourhood of the edge, and
intersecting with the neighbourhood of $R$ on which $f$ is holomorphic gives an open
neighbourhood of the edge on which $f$ is holomorphic and nonvanishing. Hence
\Cref{lem:ap_side_integral} applies on each edge and gives
\[
  \operatorname{Im} J_k = \texttt{argVariation}\ 0\ 1\ (f\circ\texttt{lineSeg}\ P_k\ P_{k+1}),
  \qquad
  \operatorname{Re} J_k = \log\lVert f(P_{k+1})\rVert-\log\lVert f(P_k)\rVert .
\]
Summing the imaginary parts over $k$ gives $\operatorname{Im}\bigl(\sum_k
J_k\bigr)=\Delta_{\partial R}\arg f$ (\Cref{def:ap_rect_argvar}). Summing the real
parts telescopes around the closed loop $P_0\to P_1\to P_2\to P_3\to P_0$:
\[
  \operatorname{Re}\Bigl(\sum_{k=0}^{3} J_k\Bigr)
  = \sum_{k=0}^{3}\bigl(\log\lVert f(P_{k+1})\rVert-\log\lVert f(P_k)\rVert\bigr) = 0
\]
(with $P_4=P_0$). Hence $\sum_k J_k$ is purely imaginary, equal to
$i\,\Delta_{\partial R}\arg f$, which is the claim.
\end{proof}

\section{The argument principle}
\label{sec:ap_main}

\begin{theorem}[Argument principle on a rectangle]
\label{thm:ap_main}
\lean{rectangle_zero_count_eq_argVariation}
\leanok
\uses{thm:ap_residue_sum, lem:ap_integral_eq_argvar, lem:ap_finite_zeros,
      def:ap_rect_argvar}
Let $U$ be open and connected, $f$ analytic on $U$ and not identically zero, and
$R=\texttt{Rectangle}\ z\ w\subseteq U$ with $z.\mathrm{re}\le w.\mathrm{re}$,
$z.\mathrm{im}\le w.\mathrm{im}$ and $f\neq0$ on $\partial R$. Then, with
$Z=\{\rho\in R:f(\rho)=0\}$,
\[
  \sum_{\rho\in Z} m_f(\rho)
  \;=\; \frac{1}{2\pi}\,\Delta_{\partial R}\arg f .
\]
That is, the number of zeros of $f$ inside $R$, counted with multiplicity, equals
$\tfrac{1}{2\pi}$ times the boundary argument variation.
\end{theorem}

\begin{proof}
\leanok
\uses{thm:ap_residue_sum, lem:ap_integral_eq_argvar}
\Cref{thm:ap_residue_sum} and \Cref{lem:ap_integral_eq_argvar} both evaluate the same
contour integral:
\[
  2\pi i \sum_{\rho\in Z} m_f(\rho)
  \;=\; \texttt{RectangleIntegral}\ (\operatorname{logDeriv} f)\ z\ w
  \;=\; i\,\Delta_{\partial R}\arg f .
\]
Cancelling the nonzero factor $i$ gives $2\pi\sum_{\rho\in Z} m_f(\rho)
=\Delta_{\partial R}\arg f$; dividing by $2\pi>0$ yields the displayed identity.
\end{proof}

\begin{remark}[Relation to \Cref{cha:n_formula}]
\label{rem:ap_scope}
\Cref{thm:ap_main} \emph{is} \Cref{thm:nf_argument_principle} of \Cref{cha:n_formula}, with
$\Delta_{\partial R}\arg f$ realised by \Cref{def:ap_rect_argvar}. In the application
of \Cref{thm:N_formula} one takes $f=\xi$ (entire and $\not\equiv0$, since
$\xi(2)\neq0$ by \Cref{lem:nf_xi_zeros}) on $U=\mathbb C$ and $R$ the rectangle with
corners $-1$ and $2+iT$; the hypothesis that $T$ is not the ordinate of a $\zeta$-zero
gives $\xi\neq0$ on $\partial R$ (\Cref{lem:nf_halving}), and the interior zero set
$Z$ is exactly the box of $\xi$-zeros counted by $N(T)$ via
\Cref{lem:nf_N_eq_xi_count}. The symmetry reduction \Cref{lem:nf_halving} then folds
$\Delta_{\partial R}\arg\xi$ to $2\,\Delta_{L_T}\arg\xi$, and the factorisation
\Cref{lem:nf_arg_poly} and \Cref{lem:nf_arg_theta} evaluates the latter, completing the
Riemann--von Mangoldt formula. Every ingredient of the present chapter reduces to the
single-pole residue calculus \Cref{thm:nf_rectangle}, mathlib's analytic-zero
codiscreteness, and \Cref{lem:nf_logderiv_residue} and
\Cref{lem:av_integral}.
\end{remark}
